We construct a cylindrical coördinate system according to right-hand orientation.
The $x$-axis will serve as the cylindrical extension of the polar $yz$-plane, centered at the hub of the wind turbine rotor $(x_0, y_0, z_0)$, downstream of the fluid flow being in the direction of increasing $x$, as is illustrated in figure \ref{fig:rotor-plane-nodes}.
For all tests we shall perform, we shall distribute points along the rotorplane in a linear manner along the radial axis for $\sfrac{r}{R} \in [\sfrac{1}{5}, 1]$, and along the azimuthal axis for $\alpha \in [0, 2\pi)$, where the ray of zero degrees coïncides with the vertical $z$-axis.
The number of points in the radial direction we choose to correspond approximately to one point per meter of blade, yielding a total number of 121 pints, regardless of the span of the interval of evaluation.
In the azimuthal direction, we place one point at every other degree, yielding a total of 180 points.
The grand total of points is thus $121 \times 180 = 21,780$.
A selection of these numerous points is shown in figure \ref{fig:evaluation_points}.
It is worth noting that the total number of panels in the \texttt{vtk}-file is 268,650, yielding the number of filaments to be 1,074,600.
These points are chosen according to a spline interpolation of the midpoint of the wind turbine blades, which is then sampled in the radial axis, and swept in the azimuthal direction to construct a faux rotorplane.

\begin{figure}[h]
  \centering
  \includegraphics[width = \textwidth]{evaluation_points.pdf}
  \caption{The evaluation rotorplane. Every third evaluation point is shown.}
  \label{fig:evaluation_points}
\end{figure}

After the evaluation of the induced velocity in the rotorplane, placing it at some specified position in the wake according to the position of the rotor hub, measured in $\sfrac{x}{D}$, the data is postprocessed in \texttt{json}-files containing the run data, and in \texttt{npz}-files containing the induced velocity.
Representing the data graphically, we have chosen two methods of displaying the induced velocity.
First is the induced velocity at $\sfrac{r}{R} = 0.8$, graphed over the azimuthal variable, illustrating the directly calculated values, with quickly calculated values of some chosen fast multipole method superimposed.
Second is the average over the $r$ axis, weighted according to the cylindrical Jacobian $r\,\dee r\,\dee\alpha$.
We shall test four such placements along the wake, namely values of $x$ corresponding to $-1D$, $-0.1D$, $0D$, and $1D$.
That is, one rotor diameter upstream of the rotorplane, onetenth rotor diameters upstream from the rotorplane, the rotorplane itself, and one rotor diameter downstream of the rotorplane.
We choose to compare the direct evaluation of the induced velocity to the \textsc{Lamb}--\textsc{Helmholtz} representation of the velocity, translated with the \textsc{rcr} method, at multipole order eight.
We shall also inspect the convergence rates of the different methods at each of these placements, measured with respect to the relative error in the velocity.
This measure of error is the standard choice in the fast multipole literature, and is calculated by taking the $\ell^2$ magnitude of the difference of the direct and quick induced velocities, divided by the $\ell^2$ magnitude of the direct induced velocity.

For the case of one rotor diameter upstream of the rotorplane, we expect the induced velocity in the azimuthal direction to be effectively zero, due to the absence of viscosity.
Inspecting the data of figures \ref{fig:neg1d_velocity_r80} and \ref{fig:neg1d_velocity_ravg}, we observe nothing unexpected.
There is indeed no induced velocity in the azimuthal direction, only relatively small but expected velocities in the radial and transversal axes.
We reckon the air must be somewhat hindered by the presence of the turbine, whence the negative value in $u^x$, and deflected around the rotorplane, whence the positive value in $u^r$.
The multipole evaluation is virtually indistinguishable, being sampled sparsely to be superimposed, so as to not visually interfere with the direct calculation.
The relative error is quite good, but not excellent at order four, as seen in figure \ref{fig:full_neg1_relative_error}, and the minimal error seems to get no smaller than on the order of $10^{-8}$.

\begin{figure}[h]
  \centering
  \includegraphics[width = .7\textwidth]{neg1d_velocity_r80.pdf}
  \caption{The induced velocity at $\sfrac{x}{D} = -1$, for $\sfrac{r}{R} = 0.8$.}
  \label{fig:neg1d_velocity_r80}
\end{figure}

\begin{figure}[h]
  \centering
  \includegraphics[width = .7\textwidth]{neg1d_velocity_ravg.pdf}
  \caption{The induced velocity at $\sfrac{x}{D} = -1$, averaged radially.}
  \label{fig:neg1d_velocity_ravg}
\end{figure}

\begin{figure}[h]
  \centering
  \includegraphics[width = .7\textwidth]{full_neg1D_relative_error.pdf}
  \caption{The relative error of the multipole induced velocity to the direct calculation at $x = -1D$.}
  \label{fig:full_neg1_relative_error}
\end{figure}

At a value of $x$ corresponding to $-0.1D$, we still do not expect to see any azimuthal induced velocity, again due to the inviscid nature of the flow.
It is clear from the plot in figures \ref{fig:negp1d_velocity_r80} and \ref{fig:negp1d_velocity_ravg}, however, that there is some induced velocity in the azimuthal direction.
We now see the three times periodic behavior expected from the presence of three distinct convected wake surfaces.
The convergence of the relative error of the multipole expansions is now improved approximately by an order of magnitude, as is clear from figure \ref{fig:full_negp1_relative_error}.
This is not unexpected, as we are now calculating more interactions directly, instead of using multipoles, when the targets are physically closer to the sources.

\begin{figure}[h]
  \centering
  \includegraphics[width = .7\textwidth]{negp1d_velocity_r80.pdf}
  \caption{The induced velocity at $\sfrac{x}{D} = -0.1$, for $\sfrac{r}{R} = 0.8$.}
  \label{fig:negp1d_velocity_r80}
\end{figure}

\begin{figure}[h]
  \centering
  \includegraphics[width = .7\textwidth]{negp1d_velocity_ravg.pdf}
  \caption{The induced velocity at $\sfrac{x}{D} = -0.1$, averaged radially.}
  \label{fig:negp1d_velocity_ravg}
\end{figure}


\begin{figure}[h]
  \centering
  \includegraphics[width = .7\textwidth]{full_negp1D_relative_error.pdf}
  \caption{The relative error of the multipole induced velocity to the direct calculation at $x = -0.1D$.}
  \label{fig:full_negp1_relative_error}
\end{figure}

For the rotorplane, the corresponding figures \ref{fig:rotorplane_velocity_r80} and \ref{fig:rotorplane_velocity_ravg} seem to show jumps in velocity for values of $\alpha$ around \SI{118}{\degree}, \SI{238}{\degree}, and \SI{358}{\degree}.
These do not correspond to the placement of the wind turbine blades, as is evident from figure \ref{fig:rotor-plane-nodes}.
Due to the orientation of that figure, the angles are inverted, as we are in effect viewing from below the cylindrical coördinate system the plots along the azimuthal angle use.
We are led to conclude that this jump is due to newly shed wakes that are now out of phase with the blades.
The relative error in the multipole accelerated induced velocity calculation is excellent, as seen in \ref{fig:full_rotorplane_relative_error}.

\begin{figure}[h]
  \centering
  \includegraphics[width = .7\textwidth]{rotorplane_velocity_r80.pdf}
  \caption{The induced velocity at the rotorplane, for $\sfrac{r}{R} = 0.8$.}
  \label{fig:rotorplane_velocity_r80}
\end{figure}

\begin{figure}[h]
  \centering
  \includegraphics[width = .7\textwidth]{rotorplane_velocity_ravg.pdf}
  \caption{The induced velocity at the rotorplane, averaged radially.}
  \label{fig:rotorplane_velocity_ravg}
\end{figure}

\begin{figure}[h]
  \centering
  \includegraphics[width = .7\textwidth]{full_rotorplane_relative_error.pdf}
  \caption{The relative error of the multipole induced velocity to the direct calculation in the rotorplane.}
  \label{fig:full_rotorplane_relative_error}
\end{figure}

Downstream at one rotor diameter from the rotorplane, the induced velocities seem to have converged to one value over the azimuthal axis for the radially averaged velocity of figure \ref{fig:1d_velocity_ravg}.
The induced velocity coïncides with the farwake modeling implemented in \texttt{3DFloat}, where panel nodes are convected according to $u^x = \SI{4}{\meter\per\second}$.
We do still see some azimuthal velocity, and there is some variation in the azimuthal direction for the induced velocities at the eighty percent rotor radius in figure \ref{fig:1d_velocity_r80}.
The multipole accuracy has again decreased, the error increasing by an order of magnitude compared to the rotorplane.

\begin{figure}[h]
  \centering
  \includegraphics[width = .7\textwidth]{1d_velocity_r80.pdf}
  \caption{The induced velocity at $\sfrac{x}{D} = 1$, for $\sfrac{r}{R} = 0.8$.}
  \label{fig:1d_velocity_r80}
\end{figure}

\begin{figure}[h]
  \centering
  \includegraphics[width = .7\textwidth]{1d_velocity_ravg.pdf}
  \caption{The induced velocity at $\sfrac{x}{D} = 1$, averaged radially.}
  \label{fig:1d_velocity_ravg}
\end{figure}

\begin{figure}[h]
  \centering
  \includegraphics[width = .7\textwidth]{full_1D_relative_error.pdf}
  \caption{The relative error of the multipole induced velocity to the direct calculation at $x = 1D$.}
  \label{fig:full_1D_relative_error}
\end{figure}

At the rotorplane we have also measured the induced velocity for an annulus $\sfrac{r}{R} \in [\sfrac{1}{5}, 2]$ in the rotorplane, averaged over the azimuthal angle, to insure that the velocity decays as expected outside the rotorplane, which is shown in figure \ref{fig:full_wide_all}.
The azimuthal induced velocity seems to have completely vanished, which is not unsurprising, considering the shape of its curve along $\alpha$ in figure \ref{fig:rotorplane_velocity_r80}.
Its magnitude is in reality of around \SI{10e-3}{\meter\per\second}, as is illustrated in figure \ref{fig:full_wide_ua}.

\begin{figure}[h]
  \centering
  \includegraphics[width = .7\textwidth]{full_wide_all.pdf}
  \caption{The average over the azimuth of all three cylindrical velocity components.}
  \label{fig:full_wide_all}
\end{figure}

\begin{figure}[h]
  \centering
  \includegraphics[width = .7\textwidth]{full_wide_ua.pdf}
  \caption{The average over the azimuth of the azimuthal velocity component.}
  \label{fig:full_wide_ua}
\end{figure}

In order to mitigate the increasingly long computational time for the direct unsteady vortex lattice method, the wake is cut at some point to reduce the total number of panels.
As they are convected, their influence at the rotorplane of course diminishes, so the removal of these panels should seemingly have little effect on the rotorplane calculations.
We see from such testing, as in figures \ref{fig:cutoff_rotorplane_average_ux} and \ref{fig:cutoff_rotorplane_average_velocity}, that the induced velocity seems to hold its shape for cuts up to fifteen rotor diameters.

\begin{figure}[h]
  \centering
  \includegraphics[width = .7\textwidth]{cutoff_rotorplane_average_ux.pdf}
  \caption{The velocity component $u^x$ in the rotorplane calculated at different cutoff distances.}
  \label{fig:cutoff_rotorplane_average_ux}
\end{figure}

\begin{figure}[h]
  \centering
  \includegraphics[width = .7\textwidth]{cutoff_rotorplane_average_velocity.pdf}
  \caption{The weighted average of the velocity in the rotorplane calculated at different cutoff distances.}
  \label{fig:cutoff_rotorplane_average_velocity}
\end{figure}
